\subsection{CAPACITABILITY OF SOUSLIN SETS}

DEFINITION 8. Let X be a Hausdorff topological space. A capacity on X is a mapping f of the set \(\mathfrak{P}(\mathbf{X})\) of all subsets of X into the extended real line \(\overline{\mathbf{R}}\) , satisfying the following conditions:

\((\mathbf{CA}_{\mathbf{I}})\) If \(\mathbf{A} \subset \mathbf{B}\) , then \(f(\mathbf{A}) \leqslant f(\mathbf{B})\) .

\((\mathbf{CA}_{\mathbf{II}})\) If \((\mathbf{A}_n)\) is any increasing sequence of subsets of \(\mathbf{X}\) , then

\[
f \left(\bigcup_ {n} \mathrm{A} _ {n}\right) = \sup _ {n} f (\mathrm{A} _ {n}).
\]

\((\mathbf{CA}_{\mathbf{III}})\) If \((\mathbf{K}_n)\) is any decreasing sequence of compact subsets of \(\mathbf{X}\) , then

\[
f \left(\bigcap_ {n} \mathrm{K} _ {n}\right) = \inf _ {n} f (\mathrm{K} _ {n}).
\]

Examples. * Let \(\mu\) be a positive measure on a locally compact space X; then the corresponding outer measure \(\mu^{*}\) is a capacity on X. It can be shown that in Euclidean space \(\mathbf{R}^{n}\) ( \(n \geqslant 3\) ), the ``Newtonian outer capacity'' is a capacity in the sense of Definition 8. *

DEFINITION 9. Let \(f\) be a capacity on \(\mathbf{X}\) ; a subset \(\mathbf{A}\) of \(\mathbf{X}\) is said to be capacitable (with respect to \(f\) ) if \(f(\mathbf{A}) = \sup_k f(\mathbf{K})\) , where \(\mathbf{K}\) runs through the set of compact subsets of \(\mathbf{A}\) .

* For example, if \(f\) is an outer measure \(\mu^{*}\) , every open set is capacitable; the capacitable sets A such that \(\mu^{*}(A) < +\infty\) are precisely the \(\mu\) -integrable sets. *

PROPOSITION 14. Let K be a compact space and let f be a capacity on K. If A is the intersection of a decreasing sequence \((\mathbf{A}_{n})\) of subsets of K, each of which is a countable union of closed sets, then A is capacitable.

It is enough to show that, for each \(a < f(A)\) , there is a closed set \(\mathbf{C} \subset \mathbf{A}\) such that \(f(\mathbf{C}) \geqslant a\) . Let us first show that there exists a sequence \((\mathbf{B}_n)_{n \geqslant 1}\) of closed sets such that \(\mathbf{B}_n \subset \mathbf{A}_n\) and such that, if we define a sequence \((\mathbf{C}_n)\) inductively by the conditions \(\mathbf{C}_0 = \mathbf{A}\) , \(\mathbf{C}_n = \mathbf{C}_{n-1} \cap \mathbf{B}_n\) for \(n \geqslant 1\) , then \(f(\mathbf{C}_n) > a\) for each \(n \geqslant 0\) . Suppose the \(\mathbf{B}_i\) have been defined for \(i < n\) ; by hypothesis we have \(\mathbf{C}_{n-1} \subset \mathbf{A} \subset \mathbf{A}_n\) and \(f(\mathbf{C}_{n-1}) > a\) ; since \(\mathbf{A}_n\) is the union of an increasing sequence of closed sets \(\mathbf{D}_j\) , it follows from \((\mathbf{CA}_{\mathbf{II}})\) that

\[
f \left(\mathrm{C} _ {n - 1}\right) = \sup _ {j} f \left(\mathrm{C} _ {n - 1} \cap \mathrm{D} _ {j}\right).
\]

Hence there is an index \(j\) such that \(f(\mathbf{C}_{n-1} \cap \mathbf{D}_j) > a\) , and we may take \(\mathbf{B}_n = \mathbf{D}_j\) .

Now let \(\mathbf{C} = \bigcap_{n}\mathbf{C}_{n}\) . Since \(\mathbf{A} = \bigcap_{n}\mathbf{A}_{n}\) and \(\mathbf{B}_n\subset \mathbf{A}_n\) we have

\[
\mathrm{C} = \bigcap_ {n} \mathrm{B} _ {n};
\]

the set C is therefore compact and contained in A.

Let \(B_n' = B_1 \cap B_2 \cap \cdots \cap B_n\) ; \((B_n')\) is a decreasing sequence of compact subsets of K; as \(C_n \subset C_i \subset B_i\) for i \textless{} n we also have \(C_n \subset B_n'\) . By \((CA_{III})\) , \(f(C) = \inf_{n} f(B_n')\) , and since \(C \subset C_n \subset B_n'\) we also have \(f(C) = \inf_{n} f(C_n) \geqslant a\) . This completes the proof.

THEOREM 5. Let X be a metrizable space and let Y be a relatively compact Souslin subspace of X. Then Y is capacitable with respect to every capacity f on X.

We have seen (no. 2, Proposition 9) that there exists a compact space K, a decreasing sequence \((\mathbf{A}_{n})\) of subsets of K, each of which is a countable union of compact sets, and a continuous mapping \(\varphi: K \to X\) such that \(Y = \varphi\left(\bigcap_{n} A_{n}\right)\) . By Proposition 14, \(\bigcap_{n} A_{n}\) is capacitable with respect to every capacity on K. Theorem 5 is therefore a consequence of the following proposition:

PROPOSITION 15. Let \(\varphi\) be a continuous mapping of a Hausdorff space K into a Hausdorff space X, and let \(f\) be a capacity on X. If for each subset A of K we put \(g(A) = f(\varphi(A))\) , then \(g\) is a capacity on K; moreover, if A is capacitable with respect to \(g\) , then \(\varphi(A)\) is capacitable with respect to \(f\) .

It is clear that \(g\) satisfies axioms \((\mathbf{CA}_{\mathbf{I}})\) and \((\mathbf{CA}_{\mathbf{II}})\) . On the other hand, let \((\mathbf{C}_n)\) be a decreasing sequence of compact subsets of \(\mathbf{K}\) , and let \(\mathbf{C} = \bigcap_{n} \mathbf{C}_n\) ; the sets \(\varphi(\mathbf{C})\) are compact, and their intersection is \(\varphi(\mathbf{C})\) ; for this intersection certainly contains \(\varphi(\mathbf{C})\) , and if \(x \in \varphi(\mathbf{C}_n)\) for all \(n\) , then the sets \(\overline{\varphi}^1(x) \cap \mathbf{C}_n\) form a decreasing sequence of non-empty compact subsets of \(\mathbf{K}\) , and therefore their intersection is not empty. We have therefore \(f(\varphi(\mathbf{C})) = \inf_{n} f(\varphi(\mathbf{C}_n))\) , that is \(g(\mathbf{C}) = \inf_{n} g(\mathbf{C}_n)\) ; thus \(g\) satisfies axiom \((\mathbf{CA}_{\mathbf{III}})\) and is therefore a capacity on \(\mathbf{K}\) .

Now let A be a subset of K which is capacitable with respect to g; if \(a < f(\varphi(A)) = g(A)\) , then there is a compact set \(C \subset A\) such that \(g(C) \geqslant a\) ; thus \(\varphi(C)\) is a compact set contained in \(\varphi(A)\) , and \(f(\varphi(C)) \geqslant a\) . This shows that \(\varphi(A)\) is capacitable with respect to f, and completes the proof of Proposition 15 and hence of Theorem 5.

Remark. * If \(\mu\) is a positive measure on a locally compact metrizable space X, then every Souslin subset A of X is \(\mu\) -measurable. For if K is any compact subset of X, K∩A is a relatively compact Souslin set, hence is capacitable with respect to \(\mu^{*}\) and consequently \(\mu\) -integrable. Note that the complement in X of a Souslin set, although not in general a Souslin set, is \(\mu\) -measurable *.

\appendix
\setcounter{chapter}{9}
\renewcommand{\thechapter}{\Roman{chapter}}
\setcounter{subsection}{0}
\renewcommand{\thesubsection}{\arabic{subsection}}

